\chapter{Complexity — Circuit Sat Reduction Spec}
%%% generated by the blueprint pipeline from TCSlib.Complexity.CircuitComplexity.CircuitSatReductionSpec
\section{Overview}

Declarations of the Lean module \texttt{TCSlib.Complexity.CircuitComplexity.CircuitSatReductionSpec}.

\section{Declarations}

\begin{definition}[Emitter phases]
\lean{BoolCircuit.CktSatReduction.Emitter.Ph}
\label{BoolCircuit.CktSatReduction.Emitter.Ph}
\leanok
\uses{BoolCircuit.GateKind}
Define the emitter’s phase set to consist of the initial count-reading phase,
gate-reading phase, first-input-reading phase, a second-input-reading phase carrying a
Boolean value, and two phases indexed by a gate kind and one of three tape indices.
\end{definition}

\begin{definition}[Local transition outcomes]
\lean{BoolCircuit.CktSatReduction.Emitter.LA}
\label{BoolCircuit.CktSatReduction.Emitter.LA}
\leanok
\uses{BoolCircuit.CktSatReduction.Emitter.Ph, BoolCircuit.GateKind}
Define a local outcome to be one of three possibilities: halt with an optional Boolean
output; continue in a specified emitter phase with an optional tape-index increment; or
begin gate emission for a specified gate kind and tape index.
\end{definition}

\begin{definition}[Local phase transition]
\lean{BoolCircuit.CktSatReduction.Emitter.lstep}
\label{BoolCircuit.CktSatReduction.Emitter.lstep}
\leanok
\uses{BoolCircuit.CircuitEval.kindOf, BoolCircuit.CktSatReduction.Emitter.LA, BoolCircuit.CktSatReduction.Emitter.Ph}
Define the local transition from an emitter phase and an optional input bit as follows.
With no input bit, it halts without output. In the initial phase, a true bit continues
in that phase and increments the tape index, while a false bit advances to gate reading;
in gate reading, true advances to reading the first gate input and false halts with
false output. The first input bit is stored in the next phase, and the second determines
the gate kind and starts the first indexed gate phase. In that phase, false begins gate
emission, while true halts without output at the last tape index and otherwise advances
to the corresponding unindexed phase. In the unindexed phase, true remains there and
increments the index; false returns to the indexed phase at the next index.
\end{definition}

\begin{definition}[Emitter instructions]
\lean{BoolCircuit.CktSatReduction.Emitter.Instr}
\label{BoolCircuit.CktSatReduction.Emitter.Instr}
\leanok
Define an instruction to be either an instruction to output a Boolean value or an
instruction to print at one of the three work-tape indices.
\end{definition}

\begin{definition}[Three-tape index reduction]
\lean{BoolCircuit.CktSatReduction.Emitter.vtx}
\label{BoolCircuit.CktSatReduction.Emitter.vtx}
\leanok
For each natural number $v$, define its corresponding tape index to be the remainder of
$v$ upon division by $3$.
\end{definition}

\begin{definition}[Literal instruction sequence]
\lean{BoolCircuit.CktSatReduction.Emitter.litProg}
\label{BoolCircuit.CktSatReduction.Emitter.litProg}
\leanok
\uses{BoolCircuit.CktSatReduction.Emitter.Instr, BoolCircuit.CktSatReduction.Emitter.vtx}
For a Boolean literal consisting of a variable and a truth-value sign, its program is
the four-instruction list that first prints the vertex associated with the variable,
then outputs true, false, and the literal’s sign, in that order.
\end{definition}

\begin{definition}[Program for a clause formula]
\lean{BoolCircuit.CktSatReduction.Emitter.progOf}
\label{BoolCircuit.CktSatReduction.Emitter.progOf}
\leanok
\uses{BoolCircuit.CktSatReduction.Emitter.Instr, BoolCircuit.CktSatReduction.Emitter.litProg}
For a conjunctive normal form formula over natural-number variables, form a program by
concatenating, for each clause in order, an output of true, the programs of its literals
in order, and an output of false.
\end{definition}

\begin{definition}[Symbolic argument lists]
\lean{BoolCircuit.CktSatReduction.Emitter.argsSym}
\label{BoolCircuit.CktSatReduction.Emitter.argsSym}
\leanok
Define a map from the three positions to lists of natural numbers by assigning the empty
list to position $0$, the list $[1]$ to position $1$, and the list $[1,2]$ to position
$2$.
\end{definition}

\begin{definition}[Program for a gate]
\lean{BoolCircuit.CktSatReduction.Emitter.prog}
\label{BoolCircuit.CktSatReduction.Emitter.prog}
\leanok
\uses{BoolCircuit.CktSatReduction.Emitter.Instr, BoolCircuit.CktSatReduction.Emitter.argsSym, BoolCircuit.CktSatReduction.Emitter.progOf, BoolCircuit.DAGGate.tseitin, BoolCircuit.GateKind}
Given a gate kind and a position $c\in\{0,1,2\}$, take the conjunctive normal form
produced by the Tseitin encoding of that gate kind at output variable $0$, with its
inputs given by the symbolic argument list for $c$, and form the program for that
formula.
\end{definition}

\begin{definition}[Execution of an instruction]
\lean{BoolCircuit.CktSatReduction.Emitter.exec}
\label{BoolCircuit.CktSatReduction.Emitter.exec}
\leanok
\uses{BoolCircuit.CktSatReduction.Emitter.Instr}
Given a vector $v:\{0,1,2\}\to\mathbb N$, an output instruction carrying a Boolean value
produces the one-element list containing that value. A print instruction at position $t$
produces a list of $v(t)$ copies of true.
\end{definition}

\begin{definition}[Bits emitted by a gate program]
\lean{BoolCircuit.CktSatReduction.Emitter.emitBits}
\label{BoolCircuit.CktSatReduction.Emitter.emitBits}
\leanok
\uses{BoolCircuit.CktSatReduction.Emitter.exec, BoolCircuit.CktSatReduction.Emitter.prog, BoolCircuit.GateKind}
Given a gate kind, a position $c\in\{0,1,2\}$, and a vector $v:\{0,1,2\}\to\mathbb N$,
the emitted bit list is obtained by executing each instruction in the program for that
gate and concatenating the resulting Boolean lists in order.
\end{definition}

\begin{definition}[Serialized formula body]
\lean{BoolCircuit.CktSatReduction.Emitter.serBody}
\label{BoolCircuit.CktSatReduction.Emitter.serBody}
\leanok
\uses{Std.Sat.CNF.serializeClause}
For a conjunctive normal form formula, form a Boolean list by concatenating, for each
clause in order, the bit true followed by the serialization of that clause.
\end{definition}

\begin{theorem}[Formula serialization decomposition]
\lean{BoolCircuit.CktSatReduction.Emitter.serialize_eq_serBody}
\label{BoolCircuit.CktSatReduction.Emitter.serialize_eq_serBody}
\leanok
\uses{BoolCircuit.CktSatReduction.Emitter.serBody, Std.Sat.CNF.serialize}
\difficulty{1}
For every CNF formula $\varphi$ over natural-number variables, its binary serialization
is the concatenation of its clause-body serialization and a final false bit.
\end{theorem}

\begin{definition}[Argument lists by position]
\lean{BoolCircuit.CktSatReduction.Emitter.argsOf}
\label{BoolCircuit.CktSatReduction.Emitter.argsOf}
\leanok
Define a map that takes a position and two natural numbers $a,b$ to the empty list when
the position is $0$, to $[a]$ when it is $1$, and to $[a,b]$ when it is $2$.
\end{definition}

\begin{theorem}[Execution of the generated program]
\lean{BoolCircuit.CktSatReduction.Emitter.flatMap_exec_progOf}
\label{BoolCircuit.CktSatReduction.Emitter.flatMap_exec_progOf}
\leanok
\uses{BoolCircuit.CktSatReduction.Emitter.exec, BoolCircuit.CktSatReduction.Emitter.litProg, BoolCircuit.CktSatReduction.Emitter.progOf, BoolCircuit.CktSatReduction.Emitter.serBody, BoolCircuit.CktSatReduction.Emitter.vtx, Std.Sat.CNF.serializeClause, Std.Sat.CNF.serializeLit}
\difficulty{5}
Let $v$ assign a natural number to each of the three indices $0,1,2$, and let $\varphi$
be a CNF formula over natural-number variables. The bit string obtained by executing the
program generated from $\varphi$ under $v$ is the body encoding of the formula obtained
by replacing each literal $(x,b)$ with $(v(x \bmod 3),b)$, preserving its polarity bit.
\end{theorem}

\begin{theorem}[Emitted bits equal serialized gate clauses]
\lean{BoolCircuit.CktSatReduction.Emitter.emitBits_eq}
\label{BoolCircuit.CktSatReduction.Emitter.emitBits_eq}
\leanok
\uses{BoolCircuit.CktSatReduction.Emitter.argsOf, BoolCircuit.CktSatReduction.Emitter.emitBits, BoolCircuit.CktSatReduction.Emitter.flatMap_exec_progOf, BoolCircuit.CktSatReduction.Emitter.prog, BoolCircuit.CktSatReduction.Emitter.serBody, BoolCircuit.DAGGate.tseitin, BoolCircuit.GateKind}
\difficulty{4}
For every gate kind $k$, arity index $c\in\{0,1,2\}$, and three-component natural-number
vector $v$, the bits emitted by the program for $(k,c)$ under $v$ equal the
serialization of the Tseitin CNF for the gate with output variable $v(0)$ and argument
list determined by $c$ from $v(1)$ and $v(2)$.
\end{theorem}

\begin{theorem}[Bound on gate program length]
\lean{BoolCircuit.CktSatReduction.Emitter.length_prog_le}
\label{BoolCircuit.CktSatReduction.Emitter.length_prog_le}
\leanok
\uses{BoolCircuit.CktSatReduction.Emitter.prog, BoolCircuit.GateKind}
\difficulty{3}
For every Boolean gate kind $k\in\{\mathrm{and},\mathrm{or},\mathrm{not}\}$ and code
$c\in\{0,1,2\}$, the instruction program encoding the corresponding gate has length at
most $56$.
\end{theorem}

\begin{definition}[Emitter state]
\lean{BoolCircuit.CktSatReduction.Emitter.ES}
\label{BoolCircuit.CktSatReduction.Emitter.ES}
\leanok
\uses{BoolCircuit.CktSatReduction.Emitter.Ph}
An emitter state consists of a phase, a current vertex $z \in \mathbb{N}$, and two
natural-number counters $a,b \in \mathbb{N}$ recording the first and second arguments
read so far.
\end{definition}

\begin{definition}[State counter vector]
\lean{BoolCircuit.CktSatReduction.Emitter.ES.vec}
\label{BoolCircuit.CktSatReduction.Emitter.ES.vec}
\leanok
\uses{BoolCircuit.CktSatReduction.Emitter.ES}
For an emitter state with counters $z$, $a$, and $b$, define its counter vector on $$
by assigning the values $z$, $a$, and $b$ to indices $0$, $1$, and $2$, respectively.
\end{definition}

\begin{definition}[Incremented emitter state]
\lean{BoolCircuit.CktSatReduction.Emitter.ES.incTo}
\label{BoolCircuit.CktSatReduction.Emitter.ES.incTo}
\leanok
\uses{BoolCircuit.CktSatReduction.Emitter.ES, BoolCircuit.CktSatReduction.Emitter.Ph}
Given an emitter state with counters $z$, $a$, and $b$, a phase $$, and an index in
$$, the incremented state has that phase and increases by one the counter selected by
the index, leaving the other two counters unchanged.
\end{definition}

\begin{theorem}[Coordinate increment of an emitter state]
\lean{BoolCircuit.CktSatReduction.Emitter.ES.vec_incTo}
\label{BoolCircuit.CktSatReduction.Emitter.ES.vec_incTo}
\leanok
\uses{BoolCircuit.CktSatReduction.Emitter.ES, BoolCircuit.CktSatReduction.Emitter.ES.incTo, BoolCircuit.CktSatReduction.Emitter.ES.vec, BoolCircuit.CktSatReduction.Emitter.Ph}
\difficulty{3}
Let an emitter state have three natural-number coordinates, written as a vector
$v=(z,a,b)$, and let $t\in\{0,1,2\}$. If the state is advanced in phase $\phi$ at
coordinate $t$, its coordinate vector is obtained from $v$ by replacing the $t$th
coordinate with $v_t+1$ and leaving the other coordinates unchanged.
\end{theorem}

\begin{definition}[Emitter run]
\lean{BoolCircuit.CktSatReduction.Emitter.emRun}
\label{BoolCircuit.CktSatReduction.Emitter.emRun}
\leanok
\uses{BoolCircuit.CktSatReduction.Emitter.ES, BoolCircuit.CktSatReduction.Emitter.ES.incTo, BoolCircuit.CktSatReduction.Emitter.ES.vec, BoolCircuit.CktSatReduction.Emitter.emitBits, BoolCircuit.CktSatReduction.Emitter.lstep}
Given an emitter state and a finite Boolean list, the emitted Boolean list is defined
recursively by processing the input from left to right according to the current phase,
advancing the state or emitting the corresponding gate bits as prescribed by the
transition rules; processing an empty input emits the empty list.
\end{definition}

\begin{definition}[Gate clauses]
\lean{BoolCircuit.CktSatReduction.Emitter.gateClauses}
\label{BoolCircuit.CktSatReduction.Emitter.gateClauses}
\leanok
\uses{BoolCircuit.CktSatReduction.Emitter.emRun}
For a finite Boolean list $w$, the gate clauses are the Boolean list obtained by running
the emitter on $w$ from the initial state, whose phase is $n$ and whose three counters
are zero.
\end{definition}

\begin{theorem}[Advancing through true input bits]
\lean{BoolCircuit.CktSatReduction.Emitter.emRun_n}
\label{BoolCircuit.CktSatReduction.Emitter.emRun_n}
\leanok
\uses{BoolCircuit.CktSatReduction.Emitter.ES.incTo, BoolCircuit.CktSatReduction.Emitter.emRun, BoolCircuit.CktSatReduction.Emitter.lstep}
\difficulty{4}
For all natural numbers $z$ and $k$ and every Boolean list $r$, running the emitter from
phase $n$ with vertex counter $z$ and both argument counters zero on a list of $k$ true
bits followed by $r$ gives the same output as running it from phase $n$ with vertex
counter $z+k$ and both argument counters zero on $r$.
\end{theorem}

\begin{theorem}[Emitter counter update in phase zero]
\lean{BoolCircuit.CktSatReduction.Emitter.emRun_u0}
\label{BoolCircuit.CktSatReduction.Emitter.emRun_u0}
\leanok
\uses{BoolCircuit.CktSatReduction.Emitter.ES.incTo, BoolCircuit.CktSatReduction.Emitter.emRun, BoolCircuit.CktSatReduction.Emitter.lstep, BoolCircuit.GateKind}
\difficulty{4}
For any gate kind, natural numbers $z,a,j$, and Boolean list $r$, processing $j$
consecutive true bits from phase $u(k,0)$ leaves the output unchanged and increases the
first argument counter from $a$ to $a+j$: the resulting run is the run from phase
$u(k,0)$ with counters $(z,a+j,0)$ on input $r$.
\end{theorem}

\begin{theorem}[Emitter counter update in phase one]
\lean{BoolCircuit.CktSatReduction.Emitter.emRun_u1}
\label{BoolCircuit.CktSatReduction.Emitter.emRun_u1}
\leanok
\uses{BoolCircuit.CktSatReduction.Emitter.ES.incTo, BoolCircuit.CktSatReduction.Emitter.emRun, BoolCircuit.CktSatReduction.Emitter.lstep, BoolCircuit.GateKind}
\difficulty{4}
For any gate kind, natural numbers $z,a,b,j$, and Boolean list $r$, processing $j$
consecutive true bits from phase $u(k,1)$ leaves the output unchanged and increases the
second argument counter from $b$ to $b+j$: the resulting run is the run from phase
$u(k,1)$ with counters $(z,a,b+j)$ on input $r$.
\end{theorem}

\begin{theorem}[Emission while reading gate arguments]
\lean{BoolCircuit.CktSatReduction.Emitter.emRun_args}
\label{BoolCircuit.CktSatReduction.Emitter.emRun_args}
\leanok
\uses{BoolCircuit.CktSatReduction.Emitter.emRun, BoolCircuit.CktSatReduction.Emitter.emRun_u0, BoolCircuit.CktSatReduction.Emitter.emRun_u1, BoolCircuit.CktSatReduction.Emitter.emitBits_eq, BoolCircuit.CktSatReduction.Emitter.lstep, BoolCircuit.CktSatReduction.Emitter.serBody, BoolCircuit.DAGGate.tseitin, BoolCircuit.GateKind, BoolCircuit.encodeList, BoolCircuit.encodeNat}
\difficulty{5}
Let $k$ be one of the gate kinds AND, OR, or NOT, let $z\in\mathbb N$, and let
$a_1,\ldots,a_m\in\mathbb N$ with $0\le m\le 2$. Encode each natural number $a$ as
$1^a0$, and encode the list $(a_1,\ldots,a_m)$ recursively by encoding the empty list as
$0$ and a nonempty list $(a,a_2,\ldots,a_m)$ as $1$ followed by the encoding of $a$ and
then the encoding of the remaining list. For every Boolean string $r$, an emitter
started in its argument-reading phase for gate kind $k$, with output index $z$ and both
argument counters zero, reads this list encoding followed by $r$ and emits the
serialized Tseitin formula for the gate with kind $k$, arguments $(a_1,\ldots,a_m)$, and
output index $z$, followed by the output of an emitter started in its gate phase with
output index $z+1$, both argument counters zero, and input $r$. Here a literal $(v,b)$
is serialized as $1^{v+1}0b$, a clause as the concatenation of its literal
serializations followed by $0$, and a formula as, for each clause in order, $1$ followed
by that clause’s serialization.
\end{theorem}

\begin{theorem}[Gate-list emission]
\lean{BoolCircuit.CktSatReduction.Emitter.emRun_gates}
\label{BoolCircuit.CktSatReduction.Emitter.emRun_gates}
\leanok
\uses{BoolCircuit.CktSatReduction.Emitter.emRun, BoolCircuit.CktSatReduction.Emitter.emRun_args, BoolCircuit.CktSatReduction.Emitter.serBody, BoolCircuit.DAGGate, BoolCircuit.DAGGate.encode, BoolCircuit.GateKind.encode, BoolCircuit.encodeList, BoolCircuit.tseitinGates}
\difficulty{5}
Let $z$ be a natural number, let $gs$ be a finite list of Boolean-circuit gates, and
suppose every gate in $gs$ has at most two arguments. For every Boolean list $r$,
starting the emitter in phase $g$ with vertex counter $z$ and both argument counters
zero, then reading the encoded list $gs$ followed by $r$, produces the serialization of
the clauses obtained by concatenating the gates’ Tseitin formulas at successive output
indices starting from $z$, followed by a false bit.
\end{theorem}

\begin{theorem}[Gate clauses encode the circuit CNF]
\lean{BoolCircuit.CktSatReduction.Emitter.gateClauses_encode}
\label{BoolCircuit.CktSatReduction.Emitter.gateClauses_encode}
\leanok
\uses{BoolCircuit.CktSatReduction.Emitter.emRun, BoolCircuit.CktSatReduction.Emitter.emRun_gates, BoolCircuit.CktSatReduction.Emitter.emRun_n, BoolCircuit.CktSatReduction.Emitter.gateClauses, BoolCircuit.CktSatReduction.Emitter.lstep, BoolCircuit.CktSatReduction.Emitter.serialize_eq_serBody, BoolCircuit.DAGCircuit, BoolCircuit.DAGCircuit.encode, BoolCircuit.encodeNat, BoolCircuit.tseitinGates, Std.Sat.CNF.serialize}
\difficulty{3}
Let $C$ be a finite acyclic Boolean circuit with natural-number input count $n$, gates
in topological order, and an output vertex among its input and gate vertices. Suppose
each gate has at most two arguments. Then running the gate-clause encoder on the circuit
encoding produces exactly the serialization of the CNF obtained by concatenating the
Tseitin clauses of the gates, starting with variable index $n$.
\end{theorem}
